\begin{reviewcomment}{6}
	On row 681 there is a formula in the text. This is not acceptable. Follow the proper format of formulas.
\end{reviewcomment}

\begin{reviewresponse}{6}
	Row 681 of the submitted manuscript belonged to the PCH results paragraph, which reported per-task statistics as inline formulas of the form $\mathrm{AUC}_{\text{mean}} = 0.561 \pm 0.053$.
	We removed inline formulas from the results prose of the PCH, SPP, and PI sections: statistics are now stated in words with the numeric values (``mean AUC $0.667$ with $53\%$ of seeds significant''), and the full figures live in the tables.
	The remaining displayed formulas are the definitions in the Design and Methodology sections and the paired leakage metric of the ablation, which is now a numbered equation.

	\textbf{Changes in the manuscript.}
	The revised result prose uses short inline numeric quantities rather than sentence-length formulas; the only new formula in this part of the evaluation is the displayed, numbered ablation metric.

	\RevisionLocation{Section ``Privacy Evaluation,'' subsection ``Policy Content Hiding,'' supervised-attack paragraph.}
	\begin{revisiontext}
		Chapter ch5 (diseases of the nervous system) meets both legs: mean AUC $0.667$ with $53\%$ of seeds significant.
		We report it as a FAIL of the criterion and examine below what the classifier separates.
	\end{revisiontext}

	\RevisionLocation{Section ``Privacy Evaluation,'' subsection ``Policy Indistinguishability,'' supervised-attack paragraph.}
	\begin{revisiontext}
		Pooled over all pairs, the AUC is $0.509\pm0.035$ with $3\%$ of seeds significant.
		Both gating pair types stay below the $0.55$ line: \texttt{semantic\_equivalent} at $0.549$ ($7\%$ significant seeds) and \texttt{weak\_equivalent} at $0.544$ ($10\%$), so PI passes.
	\end{revisiontext}

	\RevisionLocation{Section ``Privacy Evaluation,'' subsection ``Structure \& Parameter Privacy,'' supervised-attack paragraph.}
	\begin{revisiontext}
		One contrast, \texttt{nestingDepth} \texttt{medium}/\texttt{deep}, meets both legs of the criterion (mean AUC $0.563$, $53\%$ significant seeds) and is reported as FAIL.
		Its separability has the same origin as the PCH ch5 case.
	\end{revisiontext}

	\RevisionLocation{Section ``Privacy Evaluation,'' subsection ``Ablation: gas-normalization positive control,'' displayed equation defining $\Delta_{\text{leak}}$.}
	\begin{revisiontext}
		The paired leakage metric is
		\begin{equation*}
			\Delta_{\text{leak}} =
			|\mathrm{AUC}_{\mathrm{B1}} - 0.5| -
			|\mathrm{AUC}_{\mathrm{B2}} - 0.5|,
		\end{equation*}
		where positive values mean padding moved the attacker closer to chance; we bound each $\Delta_{\text{leak}}$ with a $95\%$ stability interval over $30$ paired train/test splits.
	\end{revisiontext}
\end{reviewresponse}
